% related.tex
% Focused review of the closest single-asset, deep-generative, and
% multi-asset models. Methods, validation procedures, and study results
% are described in later sections rather than repeated here.

Single-asset return models treat changing volatility, market regimes, and
extreme moves in different ways. Autoregressive conditional heteroskedasticity
(ARCH) and generalized autoregressive conditional heteroskedasticity (GARCH)
models allow conditional
variance to change over time~\cite{engle1982,bollerslev1986}. Even with
Gaussian innovations, the changing variance can produce an unconditional
return distribution with excess kurtosis. Heavy-tailed innovations can improve
the tail fit when Gaussian GARCH understates the frequency of extreme returns.
Stochastic-volatility models make volatility a latent
process~\cite{steinStein1991,heston1993,
kimShephardChib1998}. Jump-diffusion models add discontinuous price
moves~\cite{merton1976,kou2002}, but a jump component alone does not explain
persistent volatility. Markov-switching ARCH and GARCH models combine discrete
regimes with changing conditional variance~\cite{hamiltonSusmel1994,
klaassen2002}.
Hidden-state models represent returns with a small number of market regimes.
In a hidden Markov model (HMM), each latent state has a separate return distribution.
Gaussian-emission HMMs can reproduce heavy tails
and weak raw-return autocorrelation, but may still miss volatility
clustering~\cite{hamilton1989,rydenTerasvirtaAsbrink1998}. Hidden semi-Markov
models allow more flexible state durations and can improve volatility
persistence~\cite{bullaBulla2006}. An HMM does not create a jump process by
itself. Rare extreme moves must come from the state-dependent return
distributions or from a transition rule designed for such events.
Deep generative models learn return distributions and temporal features
without specifying a small set of market regimes or a parametric return
process. Generative adversarial networks (GANs) are one class of deep
generative model. Quant GANs have
reproduced several financial time-series features, including volatility
clustering and leverage effects~\cite{wieseEtAl2020}.
Sig-Wasserstein GANs use path signatures to learn features that depend on the
order of observations~\cite{niEtAl2021}. In contrast, evaluations of GAN
paths, including paths from the time-series generative adversarial network
(TimeGAN), have found too little autocorrelation in absolute
returns~\cite{takahashiEtAl2019,yoonEtAl2019timegan,kwonLee2024}. The
conflicting results show that flexible architectures do not guarantee a
match to financial stylized facts. Model design, training, and evaluation all
matter~\cite{assefaEtAl2020,stengerEtAl2024}.

Multi-asset generation also requires a model of cross-asset dependence.
Multivariate GARCH models estimate time-varying covariance matrices. The
dynamic conditional correlation model reduces the number of parameters by
separating individual volatility models from correlation
dynamics~\cite{engle2002dcc,bauwensLaurentRombouts2006}. Copulas separate each
asset's return distribution from the dependence model. Dynamic copulas allow
cross-asset dependence to change over time, while vine copulas provide flexible
high-dimensional constructions~\cite{patton2006,
aasCzadoFrigessiBakken2009}. Factor models reduce dimension by linking each
asset to a small number of shared
factors~\cite{sharpe1963,rossAPT1976,famaFrench1993}. Shared-state HMMs and
multivariate GANs can also generate joint market
moves~\cite{diasVermuntRamos2010,rizzatoEtAl2023}. A standard factor model can
resample historical residuals or fit a separate residual generator for every
asset. This paper addresses a narrower problem. We start with single-asset
generators already fitted to full returns and ask whether they can be reused in
a factor model.
The regression residuals are easy to calculate. The difficulty is that a
full-return draw already contains the asset mean and market-related variance.
Using that draw as a residual can shift the mean and count part of the variance
twice. The correction developed here removes the mean shift and variance
double count without fitting a second generator for every asset. It does not
model dependence among the remaining residuals.
